\documentclass[12pt]{article}
\usepackage[utf8]{inputenc}
\usepackage[T1]{fontenc}
\usepackage{hyperref}
\title{Classical-Middle Armenian Language Resources for K-12 Curriculum Development}
\author{Compiled by Claude Sonnet 5.5 (claude-sonnet-5-5) for Andr\'as Kornai}
\date{}
\begin{document}
\maketitle

\section{Introduction}
Classical Armenian (Grabar, ISO 639-3 xcl) is the literary language of early Armenian Christian literature. It survives as a liturgical and scholarly language. This dossier concerns the Classical stage. My search found nothing specific on Middle Armenian, so that part of the entry is not covered. Resources exist for adult and university learners, but I found no K-12 materials.

\section{Primary Educational Resources}
\begin{itemize}
\item ``Intro to Classical Armenian'': an open English introductory textbook (PDF) on krapar.org. Its preface says it is not a complete reference grammar and that it emphasizes reading. \url{https://krapar.org/docs/Intro_to_Classical_Armenian.pdf}
\item Aspirantum, ``Classical Armenian (Grabar) syllabus'': an online course, with Lesson 1 on the Armenian alphabet and the sound system of Grabar. \url{https://aspirantum.com/curriculum/classical-armenian-grabar-syllabus/armenian-alphabet-sound-system-of-grabar}
\item Leiden University course description for Classical Armenian (Dutch), which lists authors read such as Eznik and Koriwn. \url{https://studiegids.universiteitleiden.nl/courses/39198/klassiek-armeens}
\end{itemize}

\section{Linguistic and Reference Materials}
\begin{itemize}
\item \emph{New Dictionary of the Armenian Language} (Gorgias Press, 2 volumes), by three Mekhitarist fathers, with definitions in Latin. It is described as the most complete dictionary of Old Armenian. \url{https://gorgiaspress.com/new-dictionary-of-the-armenian-language-2-volume-set}
\item G\"ottingen, Armenian-related project page ``lng xcl'', with an introductory unit on Classical Armenian. \url{https://spw.uni-goettingen.de/projects/aig/lng-xcl.html}
\item Meillet, \emph{Altarmenisches Elementarbuch} (1913; reprint 1980), Jensen, \emph{Altarmenische Grammatik} (1959) and \emph{Altarmenische Chrestomathie} (1964), and Minassian, \emph{Manuel pratique d'arm\'enien ancien} (1976) were named in search results. I saw no direct URL for them.
\end{itemize}

\section{Religious and Community Resources}
Classical Armenian is the language of the early Armenian Bible translation and of church literature. I saw no specific URLs for church or community programmes in my search results.

\section{Digital and Audio Resources}
The online resources above are the only digital materials found. I found no audio recordings.

\section{Curriculum Development Context}
Classical Armenian suits a reading-focused heritage unit in Armenian-community or classical-languages programmes, with the alphabet and a few simple texts. Middle Armenian needs a separate search before it can be planned for.

\end{document}